\documentclass[11pt]{article}
\usepackage[margin=1in]{geometry}
\usepackage{amsmath,amssymb,bm,mathtools}
\usepackage{booktabs}
\usepackage{array}
\usepackage{xcolor}
\usepackage{hyperref}
\usepackage{listings}

\newcommand{\R}{\mathbf{R}}
\newcommand{\p}{\mathbf{p}}
\newcommand{\vct}{\mathbf{v}}
\newcommand{\bg}{\mathbf{b}_{\mathrm{g}}}
\newcommand{\ba}{\mathbf{b}_{\mathrm{a}}}
\newcommand{\g}{\mathbf{g}}
\newcommand{\I}{\mathbf{I}}
\newcommand{\Exp}{\operatorname{Exp}}
\newcommand{\Jr}{\mathbf{J}_r}
\newcommand{\skewop}[1]{#1^{\wedge}}
\newcommand{\atilde}{\tilde{\mathbf{a}}}
\newcommand{\omegatilde}{\tilde{\bm\omega}}

\title{ESKF State-Transition Jacobian: Continuous Error-State and Direct Discrete Derivations}
\author{}
\date{}

\begin{document}
\maketitle

\section{Purpose}
This note connects two ways of deriving the ESKF state-transition matrix $F$:
\begin{enumerate}
    \item derive the continuous-time error-state kinematics and then discretize them;
    \item start directly from the discrete nominal-state propagation and linearize it with respect to local state perturbations.
\end{enumerate}
They describe the same propagation of small errors. When the same approximation order is retained, the results agree. Direct differentiation of a discrete propagation may retain higher-order terms in $\Delta t$ that are absent from a first-order discretization $F_d \approx I+F_c\Delta t$.

\section{True, nominal, and error states}
The true state is represented by a nominal state plus a small error. For Euclidean variables,
\begin{align}
\p^t &= \p + \delta\p,\\
\vct^t &= \vct + \delta\vct,\\
\bg{}^t &= \bg + \delta\bg,\\
\ba{}^t &= \ba + \delta\ba,\\
\g^t &= \g + \delta\g.
\end{align}
For rotation, using a right perturbation,
\begin{equation}
\R^t = \R\Exp(\delta\bm\theta).
\end{equation}
The 18-dimensional error state is
\begin{equation}
\delta\bm x =
\begin{bmatrix}
\delta\p &
\delta\vct &
\delta\bm\theta &
\delta\bg &
\delta\ba &
\delta\g
\end{bmatrix}^{T}.
\end{equation}

\section{Two routes to the discrete transition matrix}
The textbook route starts from continuous-time error dynamics,
\begin{equation}
\delta\dot{\bm x}=F_c\,\delta\bm x+G_c\bm w,
\end{equation}
and then discretizes them. A first-order discretization gives
\begin{equation}
F_d \approx I+F_c\Delta t.
\end{equation}

The direct discrete route first writes the nominal propagation
\begin{equation}
\bm x_{k+1}=f_d(\bm x_k,\bm u_k),
\end{equation}
perturbs the input state, propagates it through the same discrete function, and measures the resulting output perturbation:
\begin{align}
\bm x_k^t &= \bm x_k\oplus\delta\bm x_k,\\
\bm x_{k+1}^t &= f_d(\bm x_k\oplus\delta\bm x_k,\bm u_k),\\
\bm x_{k+1}^t &= \bm x_{k+1}\oplus\delta\bm x_{k+1}.
\end{align}
Linearization gives
\begin{equation}
\delta\bm x_{k+1}\approx F_d\,\delta\bm x_k.
\end{equation}
Thus $F_d$ is the local Jacobian describing how the current error propagates into the next error.

\section{Discrete nominal-state propagation}
Using the Chapter 3 notation, the IMU measurements are denoted by
$\atilde$ and $\omegatilde$. The increments used in propagation are
\begin{align}
\bm\tau &\triangleq (\omegatilde-\bg)\Delta t.
\end{align}
The nominal propagation considered here is
\begin{align}
\p_{k+1}
&=\p_k+\vct_k\Delta t
+\frac12\R_k(\atilde-\ba)\,\Delta t^2
+\frac12\g_k\Delta t^2,\\
\vct_{k+1}
&=\vct_k+\R_k(\atilde-\ba)\,\Delta t+\g_k\Delta t,\\
\R_{k+1}
&=\R_k\Exp(\bm\tau),\\
\bg{}_{k+1}&=\bg{}_k,\\
\ba{}_{k+1}&=\ba{}_k,\\
\g_{k+1}&=\g_k.
\end{align}

\section{Position row of $F$}
Starting from
\begin{equation}
\p_{k+1}=\p_k+\vct_k\Delta t+\frac12\R_k(\atilde-\ba)\Delta t^2+\frac12\g_k\Delta t^2,
\end{equation}
we obtain immediately
\begin{align}
F_{pp}&=I,\\
F_{pv}&=I\Delta t.
\end{align}

For the rotation term, apply a right perturbation:
\begin{equation}
\R_k^t=\R_k\Exp(\delta\bm\theta_k)
\approx \R_k(I+\skewop{\delta\bm\theta_k}).
\end{equation}
Then
\begin{align}
\R_k^t(\atilde-\ba)
&\approx \R_k(\atilde-\ba)+\R_k\skewop{\delta\bm\theta_k}(\atilde-\ba)\\
&=\R_k(\atilde-\ba)-\R_k\skewop{(\atilde-\ba)}\,\delta\bm\theta_k,
\end{align}
where $\skewop{\delta\bm\theta}(\atilde-\ba)=-\skewop{(\atilde-\ba)}\delta\bm\theta$. Therefore
\begin{equation}
\boxed{F_{p\theta}=-\frac12\R_k\skewop{(\atilde-\ba)}\Delta t^2.}
\end{equation}
Because the accelerometer term is $\atilde-\ba$,
\begin{equation}
\boxed{F_{pb_a}=-\frac12\R_k\Delta t^2.}
\end{equation}
The gravity contribution is
\begin{equation}
\boxed{F_{pg}=\frac12I\Delta t^2.}
\end{equation}
There is no direct $b_g$ dependence in this discrete position update, hence
\begin{equation}
F_{pb_g}=0.
\end{equation}
The complete position error propagation is therefore
\begin{align}
\delta\p_{k+1}
={}&\delta\p_k+\Delta t\,\delta\vct_k
-\frac12\R_k\skewop{(\atilde-\ba)}\Delta t^2\delta\bm\theta_k\\
&-\frac12\R_k\Delta t^2\delta\ba{}_k
+\frac12I\Delta t^2\delta\g_k.
\end{align}

\section{Velocity row of $F$}
From
\begin{equation}
\vct_{k+1}=\vct_k+\R_k(\atilde-\ba)\Delta t+\g_k\Delta t,
\end{equation}
we obtain
\begin{align}
F_{vp}&=0,\\
F_{vv}&=I,\\
\boxed{F_{v\theta}}&=\boxed{-\R_k\skewop{(\atilde-\ba)}\Delta t},\\
F_{vb_g}&=0,\\
\boxed{F_{vb_a}}&=\boxed{-\R_k\Delta t},\\
\boxed{F_{vg}}&=\boxed{I\Delta t}.
\end{align}
Thus
\begin{equation}
\delta\vct_{k+1}
=\delta\vct_k
-\R_k\skewop{(\atilde-\ba)}\Delta t\,\delta\bm\theta_k
-\R_k\Delta t\,\delta\ba{}_k
+I\Delta t\,\delta\g_k.
\end{equation}

\section{Rotation row of $F$ and the manif Jacobians}
The rotation propagation is
\begin{equation}
\R_{k+1}=\R_k\oplus\bm\tau=\R_k\Exp(\bm\tau),
\qquad
\bm\tau=(\omegatilde-\bg)\Delta t.
\end{equation}
A \texttt{manif} call of the form
\begin{lstlisting}[language=C++,basicstyle=\ttfamily\small]
new_R = R_.rplus(tau, J_R, J_tau);
\end{lstlisting}
returns the propagated rotation together with the local Jacobians
\begin{align}
J_R &= \frac{\partial(\R\oplus\bm\tau)}{\partial\R},\\
J_\tau &= \frac{\partial(\R\oplus\bm\tau)}{\partial\bm\tau}.
\end{align}
For the right-plus convention on $SO(3)$,
\begin{align}
\boxed{J_R=\Exp(-\bm\tau)},\\
\boxed{J_\tau=J_r(\bm\tau)}.
\end{align}
Therefore
\begin{equation}
\boxed{F_{\theta\theta}=J_R=\Exp(-\bm\tau).}
\end{equation}
Since
\begin{equation}
\frac{\partial\bm\tau}{\partial\bg}=-I\Delta t,
\end{equation}
the chain rule gives
\begin{equation}
\boxed{F_{\theta b_g}=-J_\tau\Delta t=-J_r(\bm\tau)\Delta t.}
\end{equation}
Hence
\begin{equation}
\delta\bm\theta_{k+1}
=J_R\delta\bm\theta_k-J_\tau\Delta t\,\delta\bg{}_k.
\end{equation}

For a small IMU interval,
\begin{align}
J_R
&=\Exp(-\bm\tau)
\approx I-\skewop{(\omegatilde-\bg)}\Delta t+O(\Delta t^2),\\
J_r(\bm\tau)
&\approx I-\frac12\skewop{\bm\tau}+O(\|\bm\tau\|^2).
\end{align}
Consequently
\begin{equation}
-J_r(\bm\tau)\Delta t
=-I\Delta t+O(\Delta t^2),
\end{equation}
which explains the textbook first-order approximation $F_{\theta b_g}\approx-I\Delta t$.

\section{Bias and gravity rows}
With random-walk driving noise omitted from the deterministic state-transition Jacobian,
\begin{align}
\delta\bg{}_{k+1}&=\delta\bg{}_k,\\
\delta\ba{}_{k+1}&=\delta\ba{}_k,\\
\delta\g_{k+1}&=\delta\g_k.
\end{align}
Thus
\begin{equation}
F_{b_gb_g}=I,\qquad
F_{b_ab_a}=I,\qquad
F_{gg}=I,
\end{equation}
and their other deterministic blocks are zero.

\section{Complete direct-discrete state-transition matrix}
Using
\begin{equation}
\bm\tau=(\omegatilde-\bg)\Delta t,
\end{equation}
the direct Jacobian of the stated discrete nominal propagation is
\begin{equation}
\boxed{
F_d=
\begin{bmatrix}
I & I\Delta t & -\frac12\R_k\skewop{(\atilde-\ba)}\Delta t^2 & 0 & -\frac12\R_k\Delta t^2 & \frac12I\Delta t^2\\
0 & I & -\R_k\skewop{(\atilde-\ba)}\Delta t & 0 & -\R_k\Delta t & I\Delta t\\
0 & 0 & J_R & -J_\tau\Delta t & 0 & 0\\
0 & 0 & 0 & I & 0 & 0\\
0 & 0 & 0 & 0 & I & 0\\
0 & 0 & 0 & 0 & 0 & I
\end{bmatrix},
}
\end{equation}
where
\begin{equation}
J_R=\Exp(-\bm\tau),
\qquad
J_\tau=J_r(\bm\tau).
\end{equation}

\section{Why the textbook position row is simpler}
For the continuous error-state system, the position equation is
\begin{equation}
\delta\dot{\p}=\delta\vct.
\end{equation}
Therefore the continuous-time block from rotation error to position-error rate is zero:
\begin{equation}
(F_c)_{p\theta}=0.
\end{equation}
If the discrete transition is approximated only to first order,
\begin{equation}
F_d\approx I+F_c\Delta t,
\end{equation}
then
\begin{equation}
(F_d)_{p\theta}=0.
\end{equation}
The direct discrete position equation, however, contains the term
\begin{equation}
\frac12\R_k(\atilde-\ba)\Delta t^2,
\end{equation}
so its Jacobian contains
\begin{equation}
F_{p\theta}=-\frac12\R_k\skewop{(\atilde-\ba)}\Delta t^2.
\end{equation}
This is a second-order term. It is not a contradiction: it is simply discarded by the first-order approximation $I+F_c\Delta t$. A higher-order expansion
\begin{equation}
e^{F_c\Delta t}
=I+F_c\Delta t+\frac12F_c^2\Delta t^2+\cdots
\end{equation}
recovers such indirect couplings, e.g. rotation error $\rightarrow$ velocity error $\rightarrow$ position error.

Similarly,
\begin{equation}
J_R=\Exp\!\left(-(\omegatilde-\bg)\Delta t\right)
=I-\skewop{(\omegatilde-\bg)}\Delta t+O(\Delta t^2),
\end{equation}
so the exact/local discrete Lie-group Jacobian and the textbook first-order rotation block agree to first order.

\section{Main interpretation}
The essential quantity is not ``the nominal-state Jacobian'' versus ``the error-state Jacobian'' as two unrelated objects. The transition matrix $F_d$ describes how a small state error before propagation becomes a small state error after propagation:
\begin{equation}
\boxed{\delta\bm x_{k+1}\approx F_d\,\delta\bm x_k.}
\end{equation}
There are two useful derivation routes:
\begin{align*}
\text{continuous route: }&
\text{true kinematics}
\rightarrow \text{continuous error dynamics}
\rightarrow F_c
\rightarrow F_d,\\
\text{discrete route: }&
\text{discrete nominal propagation}
\rightarrow \text{local perturbation}
\rightarrow F_d.
\end{align*}
The \texttt{manif} Jacobians implement the second viewpoint naturally for Lie-group variables such as $R\in SO(3)$. For Euclidean variables such as $p$ and $v$, the same idea reduces to ordinary differentiation.

\section{Implementation consistency note}
If the discrete equations for $p_{k+1}$ and $v_{k+1}$ use the pre-propagation rotation $R_k$, their direct Jacobians should also be evaluated at $R_k$. Updating the stored rotation to $R_{k+1}$ before constructing those Jacobian blocks changes the linearization point and is not exactly the Jacobian of the discrete equations written above.

\end{document}
